\section{Introduction}\label{sec:classical}

The information-theoretic lower bound for sorting has a simple adversarial
interpretation. Keep the set of total orders consistent with the answers
given so far, and answer each comparison so as to retain a larger half of
that set. Starting from $n!$ orders, this adversary forces at least
$\log_2(n!)$ comparisons. Its direct implementation requires counting linear
extensions of a partial order, a $\#\mathrm P$-complete
problem~\cite{brightwell}. The question is whether one can obtain the same
leading constant with an efficiently computable adversary.

We consider $n$ distinct labelled keys. A sorting algorithm chooses pairs
adaptively, and must be correct for every strict total order on the labels.
An adversary is \emph{legal} if every finite sequence of its answers is
consistent with at least one such order. It sees the requested comparisons,
but need not know the sorting algorithm that generates them.
Throughout the paper, $\log$ denotes $\log_2$, and $\ln$ the natural
logarithm. The cases $n\le1$ are immediate; we assume $n\ge2$.

\begin{theorem}\label{thm:main}
There is a deterministic legal adversary, initialized using only $n$, that
performs a fixed polynomial number of bit operations per comparison and
forces every correct comparison sorter to make
\[
(1-o(1))n\log n
\]
comparisons. The term $o(1)$ tends to zero as $n\to\infty$.
\end{theorem}

Merge sort uses at most $n\ceil{\log n}$ comparisons, so the leading
coefficient is optimal. Our construction is asymptotic: its parameter
choice establishes a uniform polynomial bound, but is not intended for
practical running times.

\subsection{From binary trees to retained comparisons}

Atallah and Kosaraju~\cite{atallah} gave an efficient sorting adversary with
leading coefficient $1/2$. Related adversaries occur in the work of
Borodin, Guibas, Lynch, and Yao~\cite{borodin} on searching after partial
sorting, Richards and Vaidya~\cite{richards} on the distribution of
comparisons among keys, and Brodal, Chaudhuri, and
Radhakrishnan~\cite{brodal} on maintaining the minimum.
The binary-tree rule of Borodin et al. and Brodal et al. provides the
starting point for our construction.

Place every item at the root of an ordered binary tree. Items in disjoint
subtrees have the left-to-right order of those subtrees. A comparison changes
the placement as follows.
\begin{itemize}
\item If both items occupy the same node, move one to its left child and
the other to its right child, and answer accordingly.
\item If one item is at an ancestor of the other, move the higher item to
the child whose subtree does not contain the lower item.
\item If their nodes are unrelated by ancestry, answer from their
left-to-right order, without moving either item.
\end{itemize}
Each comparison increases the sum of item depths by at most two.
At completion, occupied nodes must be distinct and unrelated by ancestry:
otherwise two represented items can occur in either relative order, leaving
an unresolved comparison. Their binary addresses are therefore prefix-free.
If their depths are $d_1,\ldots,d_n$, the corresponding dyadic intervals
are disjoint, giving
\[
\sum_i2^{-d_i}\le1,
\qquad
2^{-\frac1n\sum_i d_i}\le\frac1n\sum_i2^{-d_i}\le\frac1n.
\]
Convexity gives the bound on $2^{-\frac1n\sum_i d_i}$. Hence $\sum_i d_i\ge n\log n$, and
at least $\tfrac12 n\log n$ comparisons are
necessary~\cite[Section~4.1]{brodal}.

The tree makes the state easy to maintain, but commits to some relations
that have never been queried. After answers $A<B$ and $C<D$ at the same
node, for example, it may place $A,C$ in the left child and $B,D$ in the
right. This also fixes $A<D$ and $C<B$. Figure~\ref{fig:retained}
contrasts this separation with retaining just the two inequalities.

\begin{figure}[tb]
\centering
\begin{tikzpicture}[>=Latex,font=\small,
  group/.style={draw=ink!65,rounded corners=2pt,fill=pale,minimum height=8mm},
  item/.style={circle,draw=ink,fill=white,minimum size=6mm,inner sep=1pt}]
\begin{scope}[shift={(-3.5,0)}]
\node at (0,1.15) {Separate at each comparison};
\node[circle,draw=ink,fill=white,minimum size=4mm] (root) at (0,.45) {};
\node[group,minimum width=19mm] (left) at (-1.25,-.55) {$A,\ C$};
\node[group,minimum width=19mm] (right) at (1.25,-.55) {$B,\ D$};
\draw[->] (root)--(left);
\draw[->] (root)--(right);
\node at (0,-1.35) {$\{A,C\}<\{B,D\}$};
\end{scope}
\begin{scope}[shift={(3.4,0)}]
\node at (0,1.15) {Retain the comparisons};
\draw[group] (-2.35,-1.03) rectangle (2.35,.78);
\node[item] (a) at (-1.25,.35) {$A$};
\node[item] (b) at (1.25,.35) {$B$};
\node[item] (c) at (-1.25,-.6) {$C$};
\node[item] (d) at (1.25,-.6) {$D$};
\draw[->,thick,accent] (a)--(b);
\draw[->,thick,accent] (c)--(d);
\draw[->,dash pattern=on 1pt off 1.5pt,thick,accent] (a)--(c);
\node at (0,-1.35) {$A<B,\quad C<D$};
\end{scope}
\end{tikzpicture}
\caption{Two answered comparisons need not determine all cross-relations.
The binary-tree rule separates the pairs spatially. A group with stored
arrows retains the freedom to answer later cross-comparisons. The dashed
arrow shows how a subsequent answer $A<C$ can be stored in the same group.}
\label{fig:retained}
\end{figure}

Kaplan, Zamir, and Zwick~\cite{kzz} improved the leading coefficient to
$37/64$ by retaining more structure before moving items down the tree.
Their constructions allow ordered pairs within a node, together with rules
for resolving interactions between these pairs. The present adversary
extends this approach: a node retains a directed graph of comparisons, and
an item moves to a child only when the graph must be simplified or its
coordinate must be prepared for a query.

\subsection{The construction and its analysis}

For a partial order $P$, the region of unit coordinates satisfying its
inequalities has volume $e(P)/n!$, where $e(P)$ is the number of linear
extensions: every coordinate order occupies volume $1/n!$~\cite{stanley}.
Thus retaining a larger-volume side would implement the ideal adversary
if we kept all items at the root and never simplified its constraints.

In our representation each group has auxiliary real positions in $(0,1)$,
constrained by its stored comparisons. An arrow $x\to y$ means $X_x<X_y$. The adversary chooses
an approximately larger-volume answer, retaining at least $1/2-\eta$ of
the group's feasible volume. Bounded degree and bounded directed-path
length allow this probability to be estimated from a neighborhood whose
size depends only on an accuracy parameter. The group itself may contain
arbitrarily many items.

To simplify a group, fix an item's position at an interior point $c$, replace
its incident arrows by interval bounds, and give it a fresh coordinate in
a child at $c$. This preserves every answered comparison. If $f_x(c)$ is
the density of the fixed coordinate, the product $Z$ of all group volumes
is multiplied by exactly $f_x(c)$. A well-chosen position can therefore
recover volume while removing constraints from the graph.

The proof assigns a small interval to each child and embeds the entire
recursive state in the unit cube. The widths of these intervals contribute
a factor $\mathcal W$. Initially $n!Z\mathcal W=n!$; when a sorter finishes,
the embedding lies in one order region and $n!Z\mathcal W\le2^n$.
We then use the potential
\[
\Psi=n!Z^\beta\mathcal W\,2^{-\tau e},
\]
where $e$ is the number of stored arrows, $\beta=1+o(1)$, and $\tau=o(1)$.
The extra volume factor and the charge for stored arrows make every complete
repair nondecreasing in $\Psi$. Each comparison reduces $\log\Psi$ by at
most $1+o(1)$, giving the theorem.

Section~\ref{sec:algorithm} specifies the algorithm before its analysis.
Sections~\ref{sec:legal} and~\ref{sec:oracle} establish legality and the
local numerical guarantees. Sections~\ref{sec:accounting}
and~\ref{sec:potential} give the packing and potential arguments, and
Section~\ref{sec:runtime} proves the uniform polynomial bound.
The proofs of the density bounds, locality, and exact local volume calculation
appear in the appendices. The manuscript source is available online.\footnote{\url{https://github.com/nicknash/sorting-adversary-proofs/tree/main/for-humans}}
